% Mobile and Embedded Computing, Lab 3. Compile: latexmk -xelatex lab03.tex
\documentclass[10pt,aspectratio=169,t]{beamer}
\usetheme[lab]{upbminimal}

\course{Mobile and Embedded Computing}
\decklabel{Lab 3}
\title{Bluetooth Low Energy}
\subtitle{A peripheral in Go and a Flutter central}
\author{Dinu-Ștefan Rusu}
\institute{FILS, Universitatea Politehnica București}
\date{Week 6}

\begin{document}

\titleframe

\begin{frame}[eyebrow=Today]{What this lab is for}
  \vskip 2mm
  \begin{grid}[2]
    \cell[\faMicrochip]{Run a peripheral}{Advertise one service with one notifying characteristic from a Go program, on an nRF52840 or ESP32-C3 board or on a Linux or Windows laptop.}
    \cell[\faBluetoothB]{Write the central}{Scan, connect, subscribe and show the value live in Flutter, and keep working after the peripheral disappears.}
  \end{grid}
  \vskip 6mm
  \taskmeta{Time}{2 hours}{Submit}{Branch \code{lab-3} in the team repository, pushed before the next lab}
\end{frame}

\begin{frame}[fragile,eyebrow=Setup]{Where your code goes}
  \begin{steps}
    \item Branch the team repository and add the two packages:
\begin{shell}
git checkout -b lab-3
flutter pub add flutter_blue_plus permission_handler
\end{shell}
    \item The peripheral lives in \code{peripheral/main.go}, in its own Go module.
    \item The app screen lives in \code{lib/lab3/ble\_screen.dart}, reached from your home screen.
  \end{steps}
  \begin{warning}
    Neither the Android emulator nor the iOS simulator has a Bluetooth radio. Everything in this lab runs on a physical phone, over a cable.
  \end{warning}
\end{frame}

\begin{frame}[embedded,fragile,eyebrow=Peripheral]{Pick a peripheral, then start it}
  \begin{columns}[T]
    \begin{column}{0.52\textwidth}
\begin{shell}
# Laptop: Linux or Windows
cd peripheral
go mod init lab3
go get tinygo.org/x/bluetooth
go run .

# Board: nRF52840, same source
tinygo flash -target=pca10056-s140v7 .
\end{shell}
    \end{column}
    \begin{column}{0.44\textwidth}
      \lead{One source file, two ways to run it}{The same program runs under \code{go run} on a laptop and on a board after \code{tinygo flash}. The Flutter side cannot tell the difference.}
      \vskip 2mm
      \begin{warning}
        The ESP32-C3 and ESP32-S3 work through espradio; a plain ESP32 needs an HCI co-processor. macOS is central-only, so a Mac cannot host the peripheral.
      \end{warning}
    \end{column}
  \end{columns}
\end{frame}

\begin{frame}[embedded,fragile,eyebrow=Reference]{The peripheral, part 1: advertise}
\begin{golang}
package main
import "time"
import "tinygo.org/x/bluetooth"
var adapter = bluetooth.DefaultAdapter
var value bluetooth.Characteristic
var svc, _ = bluetooth.ParseUUID("0000fff0-0000-1000-8000-00805f9b34fb")
var chr, _ = bluetooth.ParseUUID("0000fff1-0000-1000-8000-00805f9b34fb")

func main() {
    must(adapter.Enable())   // must panics on a non-nil error
    adv := adapter.DefaultAdvertisement()
    must(adv.Configure(bluetooth.AdvertisementOptions{
        LocalName: "MEC Lab 3", ServiceUUIDs: []bluetooth.UUID{svc}}))
\end{golang}
\end{frame}

\begin{frame}[embedded,fragile,eyebrow=Reference]{The peripheral, part 2: notify}
\begin{golang}
    must(adv.Start())
    must(adapter.AddService(&bluetooth.Service{UUID: svc,
        Characteristics: []bluetooth.CharacteristicConfig{{
            Handle: &value, UUID: chr,
            Flags: bluetooth.CharacteristicReadPermission |
                bluetooth.CharacteristicNotifyPermission}}}))
    go func() {                  // writing the handle notifies
        for i := byte(0); ; i++ {
            value.Write([]byte{i}); time.Sleep(time.Second)
        }
    }()
    select {}
}
\end{golang}
\end{frame}

\begin{frame}[eyebrow=Task I]{Ask for the radio, then list what you find}
  \begin{columns}[T]
    \begin{column}{0.55\textwidth}
      \begin{steps}
        \item Declare the Android permissions in the manifest and the iOS usage string in Info.plist.
        \item Request BLUETOOTH\_SCAN and BLUETOOTH\_CONNECT at runtime with permission\_handler.
        \item Wait for the adapter state to be on, then scan with flutter\_blue\_plus.
        \item Show one row per device: name and RSSI, sorted by RSSI.
      \end{steps}
    \end{column}
    \begin{column}{0.41\textwidth}
      \kv{Done when}{}
      \begin{checklist}
        \item A clean install shows the system dialog before any scan starts.
        \item \code{MEC Lab 3} appears in the list within five seconds of opening the screen.
        \item The RSSI number changes when you walk away from the peripheral.
        \item Denying the permission shows a message, not a crash.
      \end{checklist}
    \end{column}
  \end{columns}
\end{frame}

\begin{frame}[fragile,eyebrow=Reference]{The two files no Dart code can replace}
\begin{codeblock}
<!-- android/app/src/main/AndroidManifest.xml, inside <manifest> -->
<uses-permission android:name="android.permission.BLUETOOTH_SCAN"
                 android:usesPermissionFlags="neverForLocation"/>
<uses-permission android:name="android.permission.BLUETOOTH_CONNECT"/>
<uses-permission android:name="android.permission.ACCESS_FINE_LOCATION"
                 android:maxSdkVersion="30"/>

<!-- ios/Runner/Info.plist, inside the top-level <dict> -->
<key>NSBluetoothAlwaysUsageDescription</key>
<string>Shows the live reading from your lab peripheral.</string>
\end{codeblock}
  \vskip 2mm
  \muted{An undeclared Android permission is denied with no dialog. A missing iOS usage string terminates the app the moment it touches the Bluetooth API.}
\end{frame}

\begin{frame}[fragile,eyebrow=Reference]{Request, wait for the adapter, scan}
\begin{dart}
if (Platform.isAndroid) {
  await [Permission.bluetoothScan,
         Permission.bluetoothConnect].request();
}
await FlutterBluePlus.adapterState
    .where((s) => s == BluetoothAdapterState.on).first;

final sub = FlutterBluePlus.onScanResults
    .listen((r) => setState(() => _found = r));
FlutterBluePlus.cancelWhenScanComplete(sub);
await FlutterBluePlus.startScan(
    withServices: [Guid(kServiceUuid)],   // drop once, to debug
    timeout: const Duration(seconds: 15));
\end{dart}
\end{frame}

\begin{frame}[eyebrow=Task II]{Connect, discover, and show the value}
  \begin{columns}[T]
    \begin{column}{0.55\textwidth}
      \begin{steps}
        \item Tapping a row stops the scan and connects, with a ten-second timeout.
        \item Discover the services and find the characteristic by its UUID.
        \item Listen to the value stream, then call \code{setNotifyValue(true)}.
        \item Show the decoded counter and the time of the last update.
      \end{steps}
    \end{column}
    \begin{column}{0.41\textwidth}
      \kv{Done when}{}
      \begin{checklist}
        \item The number on screen increases by one every second.
        \item No two updates in a row show the same number.
        \item Leaving the screen stops notifications and disconnects.
        \item The UUIDs appear once, as constants, not typed in two places.
      \end{checklist}
    \end{column}
  \end{columns}
  \begin{takeaway}{Listen first, then enable.}
    Enabling notifications before the listener exists loses what arrives in between.
  \end{takeaway}
\end{frame}

\begin{frame}[fragile,eyebrow=Reference]{Find the characteristic and listen}
\begin{dart}
await device.connect(timeout: const Duration(seconds: 10));
final services = await device.discoverServices();

final c = services
    .firstWhere((s) => s.uuid == Guid(kServiceUuid))
    .characteristics
    .firstWhere((x) => x.uuid == Guid(kCharUuid));

final sub = c.onValueReceived.listen(
    (v) => setState(() => _value = v.isEmpty ? null : v.first));
device.cancelWhenDisconnected(sub);

await c.setNotifyValue(true);   // writes the descriptor for you
\end{dart}
\end{frame}

\begin{frame}[eyebrow=Task III]{Survive a disconnect without a restart}
  \begin{columns}[T]
    \begin{column}{0.55\textwidth}
      \begin{steps}
        \item Subscribe to \code{connectionState} before the first connect.
        \item On a disconnect, clear the value and show \emph{disconnected}.
        \item Retry after a two-second delay, and keep retrying while the screen is open.
        \item Rediscover the services after every reconnection.
      \end{steps}
    \end{column}
    \begin{column}{0.41\textwidth}
      \kv{Done when}{}
      \begin{checklist}
        \item Killing the peripheral turns the screen to \emph{disconnected} within five seconds.
        \item No stale number is left on screen after the disconnect.
        \item Restarting the peripheral resumes updates without restarting the app.
        \item The retry survives three kill and restart cycles in a row.
      \end{checklist}
    \end{column}
  \end{columns}
\end{frame}

\begin{frame}[fragile,eyebrow=Reference]{A reconnect loop that stops when you leave}
\begin{dart}
device.connectionState.listen((s) async {
  if (s != BluetoothConnectionState.disconnected) return;
  setState(() { _value = null; _online = false; });
  await Future.delayed(const Duration(seconds: 2));
  if (mounted && _wanted) _tryConnect();  // _wanted: false in dispose
});

Future<void> _tryConnect() async {
  try {
    await device.connect(timeout: const Duration(seconds: 10));
    await _subscribe();       // rediscover: old objects are dead
  } catch (_) { /* the listener above schedules the next try */ }
}
\end{dart}
\end{frame}

\begin{frame}[eyebrow=Troubleshooting]{Scanning and connecting}
  \vskip 2mm
  \trouble{The device list stays empty}{Permission denied, or location is off on Android 11 and earlier. Grant Nearby devices, switch location on.}
  \trouble{connect, android-code: 133}{The peripheral is not advertising, or another central still holds it. Restart it, wait two seconds, retry.}
  \trouble{The app quits on the first scan on iOS}{NSBluetoothAlwaysUsageDescription is missing. Add it, then rebuild: a hot restart does not reread Info.plist.}
  \trouble{Every row in the list is unnamed}{The advertisement carries no local name, or you read the wrong field. Use platformName.}
\end{frame}

\begin{frame}[eyebrow=Troubleshooting]{Nothing arrives}
  \vskip 2mm
  \trouble{Bad state: No element}{A UUID in the app differs from the one in the Go file, or you reused services from the previous connection.}
  \trouble{Connected, no value arrives}{setNotifyValue(true) was never called, or the listener was added after it.}
  \trouble{The counter moves, the phone is idle}{A notification only reaches a subscribed central, and the characteristic needs the notify flag.}
  \trouble{One value arrives, then nothing}{The subscription died with the connection. Discover and subscribe again inside the reconnect path.}
\end{frame}

\begin{frame}[embedded,eyebrow=Troubleshooting]{The peripheral side}
  \trouble{Bluetooth adapter not found}{On Linux the service is down or the radio is blocked. Start bluetoothd, unblock with rfkill, run again.}
  \trouble{It runs, nothing scans it}{macOS is central-only in this package. Use Linux, Windows, or the board.}
  \trouble{tinygo flash: unknown target}{Use your board's target: pca10056-s140v7 (nRF52840 DK) or xiao-esp32c3.}
  \trouble{Nothing works on an ESP32}{Only the ESP32-C3 and ESP32-S3 advertise on their own, through espradio. A plain ESP32 needs an HCI co-processor: use the nRF52840 or a laptop instead.}
\end{frame}

\begin{frame}[eyebrow=Submission]{What the grader will do}
  \vskip 2mm
  \begin{columns}[T]
    \begin{column}{0.52\textwidth}
      \kv{Checked on your branch}{}
      \begin{checklist}
        \item \code{peripheral/main.go} builds and advertises.
        \item The manifest and Info.plist entries are committed.
        \item The scan list shows name and RSSI.
        \item The live value updates once per second.
        \item The peripheral is killed and restarted, and the app recovers.
      \end{checklist}
    \end{column}
    \begin{column}{0.44\textwidth}
      \kv{How to submit}{}
      \muted{Branch \code{lab-3} in the team repository, one commit per task, pushed before the next lab. Put your two UUIDs and the byte encoding in the README, and three screenshots in the pull request.}
      \vskip 2mm
      \begin{hint}[Before you leave]
        Run the whole path once from a fresh install of the app, with the phone unplugged from the laptop.
      \end{hint}
    \end{column}
  \end{columns}
\end{frame}

\closingframe{Push before you leave.}{Branch lab-3 · peripheral, scan, live value, reconnect}

\end{document}
